\documentclass{article}

% packages
\usepackage{amsmath, amsthm, thmtools, amsfonts, amssymb, luacode, catchfile, tikzducks, hyperref, ifthen}
\ifcsname c@kobocompile\endcsname
	\usepackage[a5paper, total={1072pt, 1448pt}, margin=10pt, includeheadfoot]{geometry} % set page margins
\else
	\usepackage[a4paper, margin=50pt, includeheadfoot]{geometry}
\fi
\usepackage[shortlabels]{enumitem}
\usepackage[skip=3pt, indent=0pt]{parskip}

% language
\usepackage[bidi=basic, layout=tabular, provide=*]{babel}
\ifcsname c@english\endcsname
	\babelprovide[main, import]{english}
\else
	\babelprovide[main, import]{hebrew}
	\babelprovide{rl}
\fi
%\babelfont{rm}{Libertinus Serif}
\babelfont{rm}[Renderer=Harfbuzz]{Libertinus Serif}
\babelfont{sf}{Libertinus Sans}
\babelfont{tt}{Libertinus Mono}

% style
\AddToHook{cmd/section/before}{\clearpage}	% Add line break before section
\linespread{1.3}
\setcounter{secnumdepth}{0}		% Remove default number tags from sections, this won't do well with theorems
\AtBeginDocument{\setlength{\belowdisplayskip}{3pt}}
\AtBeginDocument{\setlength{\abovedisplayskip}{3pt}}
\graphicspath{ {../images/} }

% operators
\DeclareMathOperator\cis{cis}
\DeclareMathOperator\Sp{Sp}
\DeclareMathOperator\tr{tr}
\DeclareMathOperator\im{Im}
\DeclareMathOperator\re{Re}
\DeclareMathOperator\diag{diag}
\DeclareMathOperator*\lowlim{\underline{lim}}
\DeclareMathOperator*\uplim{\overline{lim}}
\DeclareMathOperator\rng{rng}
\DeclareMathOperator\Sym{Sym}
\DeclareMathOperator\Arg{Arg}
\DeclareMathOperator\Log{Log}
\DeclareMathOperator\dom{dom}
\DeclareMathOperator\supp{Supp}
\DeclareMathOperator\var{Var}
\DeclareMathOperator\cov{Cov}

% commands
%\renewcommand\qedsymbol{\textbf{מש''ל}}
%\renewcommand\qedsymbol{\fbox{\emoji{lizard}}}
\newcommand{\Aa}[0]{\mathcal{A}}
\newcommand{\Bb}[0]{\mathcal{B}}
\newcommand{\CC}[0]{\mathbb{C}}
\newcommand{\Cc}[0]{\mathcal{C}}
\newcommand{\EE}[0]{\mathbb{E}}
\newcommand{\FF}[0]{\mathbb{F}}
\newcommand{\Ff}[0]{\mathcal{F}}
\newcommand{\Ii}[0]{\mathcal{I}}
\newcommand{\Gg}[0]{\mathcal{G}}
\newcommand{\Ll}[0]{\mathcal{L}}
\newcommand{\Mm}[0]{\mathcal{M}}
\newcommand{\NN}[0]{\mathbb{N}}
\newcommand{\Nn}[0]{\mathcal{N}}
\newcommand{\PP}[0]{\mathbb{P}}
\newcommand{\Pp}[0]{\mathcal{P}}
\newcommand{\QQ}[0]{\mathbb{Q}}
\newcommand{\RR}[0]{\mathbb{R}}
\newcommand{\Rr}[0]{\mathcal{R}}
\newcommand{\Ss}[0]{\mathcal{S}}
\newcommand{\TT}[0]{\mathbb{T}}
\newcommand{\Uu}[0]{\mathcal{U}}
\newcommand{\Vv}[0]{\mathcal{V}}
\newcommand{\Ww}[0]{\mathcal{W}}
\newcommand{\ZZ}[0]{\mathbb{Z}}
\newcommand{\acts}[0]{\circlearrowright}
\newcommand{\explain}[2] {
	\begin{flalign*}
		 && \text{#2} && \text{#1}
	\end{flalign*}
}
\newcommand{\maketitleprint}[0]{ \begin{center}
	%\begin{tikzpicture}[scale=3]
	%	\duck[graduate=gray!20!black, tassel=red!70!black]
	%\end{tikzpicture}	
	\includegraphics[width=6cm]{cover}
\end{center}
}

% theorem commands
\newtheoremstyle{c_remark}
	{}	% Space above
	{}	% Space below
	{}% Body font
	{}	% Indent amount
	{\bfseries}	% Theorem head font
	{}	% Punctuation after theorem head
	{.5em}	% Space after theorem head
	{\thmname{#1}\thmnumber{ #2}\thmnote{ \normalfont{\text{(#3)}}}}	% head content
\newtheoremstyle{c_definition}
	{3pt}	% Space above
	{3pt}	% Space below
	{}% Body font
	{}	% Indent amount
	{\bfseries}	% Theorem head font
	{}	% Punctuation after theorem head
	{.5em}	% Space after theorem head
	{\thmname{#1}\thmnumber{ #2}\thmnote{ \normalfont{\text{(#3)}}}}	% head content
\newtheoremstyle{c_plain}
	{3pt}	% Space above
	{3pt}	% Space below
	{\itshape}% Body font
	{}	% Indent amount
	{\bfseries}	% Theorem head font
	{}	% Punctuation after theorem head
	{.5em}	% Space after theorem head
	{\thmname{#1}\thmnumber{ #2}\thmnote{ \text{(#3)}}}	% head content

\ifcsname c@english\endcsname
	\theoremstyle{plain}
	\newtheorem{theorem}{Theorem}[section]
	\newtheorem{lemma}[theorem]{Lemma}
	\newtheorem{proposition}[theorem]{Proposition}
	\newtheorem*{proposition*}{Proposition}
	%\newtheorem{corollary}[theorem]{אין חלופה עברית}

	\theoremstyle{definition}
	\newtheorem{definition}[theorem]{Definition}
	\newtheorem*{definition*}{Definition}
	\newtheorem{example}{Example}[section]
	\newtheorem{exercise}{Exercise}[section]

	\theoremstyle{remark}
	\newtheorem*{remark}{Remark}
	\newtheorem*{solution}{Solution}
	\newtheorem{conclusion}[theorem]{Conclusion}
	\newtheorem{notation}[theorem]{Notation}
\else
	\theoremstyle{c_plain}
	\newtheorem{theorem}{משפט}[section]
	\newtheorem{lemma}[theorem]{למה}
	\newtheorem{proposition}[theorem]{טענה}
	\newtheorem*{proposition*}{טענה}
	%\newtheorem{corollary}[theorem]{אין חלופה עברית}

	\theoremstyle{c_definition}
	\newtheorem{definition}[theorem]{הגדרה}
	\newtheorem*{definition*}{הגדרה}
	\newtheorem{example}{דוגמה}[section]
	\newtheorem{exercise}{תרגיל}[section]

	\theoremstyle{c_remark}
	\newtheorem*{remark}{הערה}
	\newtheorem*{solution}{פתרון}
	\newtheorem{conclusion}[theorem]{מסקנה}
	\newtheorem{notation}[theorem]{סימון}
\fi

% Questions related commands
\newcounter{question}
\setcounter{question}{1}
\newcounter{sub_question}
\setcounter{sub_question}{1}

\ifcsname c@english\endcsname
	\newcommand{\question}[1][0]{
		\ifthenelse{#1 = 0}{}{\setcounter{question}{#1}}
		\section{Question \arabic{question}}
		\addtocounter{question}{1}
		\setcounter{sub_question}{1}
	}

	\newcommand{\subquestion}[1][0]{
		\ifthenelse{#1 = 0}{}{\setcounter{sub_question}{#1}}
		\subsection{Part \alph{sub_question}}
		\addtocounter{sub_question}{1}
	}
\else
	\newcommand{\question}[1][0]{
		\ifthenelse{#1 = 0}{}{\setcounter{question}{#1}}
		\section{שאלה \arabic{question}}
		\addtocounter{question}{1}
		\setcounter{sub_question}{1}
	}

	\newcommand{\subquestion}[1][0]{
		\ifthenelse{#1 = 0}{}{\setcounter{sub_question}{#1}}
		\subsection{סעיף \localecounter{letters.gershayim}{sub_question}}
		\addtocounter{sub_question}{1}
	}
\fi

% import lua and start of document
\directlua{common = require ('../common')}

\GetEnv{AUTHOR}

% headers
\author{\AUTHOR}
\date\today

\title{פתרון מטלה 10 --- תורת ההסתברות (1), 80420}
% chktex-file 44

\DeclareMathOperator{\Supp}{Supp}

\begin{document}
\maketitle
\maketitleprint{}

\question{}
יהי $X$ משתנה מקרי שהמומנט ה־$k$ שלו קיים וסופי. \\
נוכיח שלכל $c > 0$ מתקיים
\[
	\PP(|X - \EE(X)| \ge c)
	\le \frac{\EE({\left\lvert X - \EE(X)\right\rvert}^k)}{c^k}
\]
\begin{proof}
	נבחין כי הקבוצות הבאות שקולות,
	\[
		\{ |X - \EE(X)| \ge c \}
		= \{ {|X - \EE(X)|}^k \ge c^k \}
	\]
	ולכן מאי־שוויון מרקוב והעובדה שנתון כי המומנט ה־$k$ של $X$ מוגדר,
	\[
		\PP(|X - \EE(X)| \ge c)
		= \PP({|X - \EE(X)|}^k \ge c^k)
		\le \frac{\EE({|X - \EE(X)|}^k)}{c^k}
	\]
\end{proof}

\question{}
נחשב את הפונקציה יוצרת המומנטים של $U([k])$ ונקבע את תחום הגדרתה.
\begin{solution}
	נניח ש־$X \sim U([k])$ ולכן עלינו למצוא את $M_X(t)$ לכל $t \in \RR$ או להראות ש־$M_X(t)$ לא מוגדר.
	\begin{align*}
		M_X(t)
		& = \EE(e^{t X}) \\
		& = \sum_{s \in \supp e^{t X}} s \PP(e^{t X} = s) \\
		& = \sum_{s \in \{ e^{tn} \mid n \in [k]\}} s \PP(e^{t X} = s) \\
		& = \sum_{n = 1}^{k} e^{tn} \PP(e^{tX} = e^{tn}) \\
		& = \sum_{n = 1}^{k} e^{tn} \frac{1}{k} \\
		& = \frac{1}{k} \cdot e^t \frac{e^{tk} - 1}{e^t - 1} \\
		& = \frac{e^{t(k + 1)} - e^t}{k(e^t - 1)}
	\end{align*}
	ולכן $M_X(t) = \frac{e^{t(k + 1)} - e^t}{k(e^t - 1)}$, אבל מצאנו בחישוב שלכל $t \ne 0$ הביטוי מוגדר, ולכן $M_X(t) : \ZZ \setminus \{ 0 \} \to \RR$.
\end{solution}

\question{}
יהי $X$ משתנה מקרי עם פונקציה יוצרת מומנטים $M_X(t) = at + b$ עבור $a, b \in \RR$. \\
נמצא את ערכי $a, b$ ואת התפלגות $X$.
\begin{solution}
	נתחיל בהצבה,
	\[
		M_X(0)
		= \EE(e^{0X})
		= \EE(1)
		= 1
		= a \cdot 0 + b
	\]
	ולכן $b = 1$.
	עוד אנו יודעים ש־$M_X'(0) = \EE(X)$ אבל $M_X'(t) = a$ ולכן $\EE(X) = a$, וכן $\EE(X^2) = M_X''(0) = 0$ ולכן
	\[
		\var(X)
		= \EE(X^2) - {\EE(X)}^2
		= 0 - a^2
	\]
	ומאי־שליליות השונות נובע ש־$a = 0$ בלבד, כלומר $M_X(t) = 1$ וכן $\EE(X) = 0, \var(X) = 0$, כלומר $X = 0$.
\end{solution}

\question{}
יהי $X \sim Poi(\lambda)$,
נראה שלכל $c > 1$ מתקיים
\[
	\PP(X \ge c \lambda)
	\le \frac{e^{c\lambda - \lambda}}{c^{c\lambda}}
\]
\begin{proof}
	בהרצאה מצאנו שמתקיים $M_X(t) = e^{\lambda(e^t - 1)}$, לכן מאי־שוויון צ'רנוף לכל $t > 0$,
	\[
		\PP(X \ge c \lambda)
		\le \frac{M_X(t)}{e^{t c \lambda}}
		= \frac{e^{c\lambda (e^t - 1)}}{e^{t c \lambda}}
	\]
	לכן בפרט כאשר $t = \log c$ מתקיים,
	\[
		\PP(X \ge c \lambda)
		\le \frac{e^{\lambda (e^{\log c} - 1)}}{e^{\log(c) c \lambda}}
		= \frac{e^{\lambda c - \lambda}}{c^{c \lambda}}
	\]
	כפי שרצינו.
\end{proof}

\question{}
יהי $n \in \NN$ ונגדיר טבלה בגודל $n \times n$, נעבור על כל צלע בטבלה ונמחק אותה בהסתברות $\frac{1}{2}$ באופן בלתי־תלוי. \\
יהי $X$ מספר הריבועים שלא נמחקה אף צלע שלהם, נוכיח על־ידי שימוש באי־שוויון הופדינג שלכל $c > 0$,
\[
	\PP(X \ge \frac{n^2}{16} + cn) \le 2 e^{-\frac{c^2}{4}}
\]
\begin{proof}
	נגדיר $Y_{i, j} \sim Ber(\frac{1}{2})$ האם הצלע ה־$i, j$ קיימת (מטעמי סימטריה מותר לנו לבחור כן). \\
	נגדיר גם $X_{i, j} = Y_{i, j} \cdot Y_{i + 1, j} \cdot Y_{i, j + 1} \cdot Y_{i + 1, j + 1} \sim Ber(\frac{1}{16})$ המייצג שהריבוע ה־$i, j$ קיים. \\
	נשים לב ש־$|X_{i, j}| \le 1$ מההגדרה ולכן יהיה לנו תוקף לאי־שוויון הופדינג בהמשך. \\
	לבסוף גם $X = \sum_{1 \le i, j \le n} X_{i, j}$. \\
	נראה דוגמה למבנה של טבלה כזו,
	\begin{center}
		\begin{tabular}{|c | c | c |}
			\hline
			A & B & A \\
			\hline
			B & A & B \\
			\hline
			A & B & A \\
			\hline
		\end{tabular}
	\end{center}
	כדי להשתמש באי־שוויון הופדינג נחלק את המשבצות לקבוצות בלתי־תלויות, נקרא למשבצת $X_{i, j}$ זוגית אם ורק אם $i + i \equiv 0 \mod 2$, ונבחין שהקבוצה $S_0$ קבוצת המשבצות הזוגיות ו־$S_1$ קבוצת המשבצות האי־זוגיות בלתי־תלויות.
	נגדיר פורמלית,
	\[
		S_0 = \sum_{\substack{1 \le i, j \le n \\ i + j \equiv 0 \mod 2}} X_{i, j},
		\qquad
		S_1 = \sum_{\substack{1 \le i, j \le n \\ i + j \equiv 1 \mod 2}} X_{i, j}
	\]
	וכן $X = S_0 + S_1$.
	נבחן את $S_0$,
	\[
		\EE(S_0)
		= \sum_{\substack{1 \le i, j \le n \\ i + j \equiv 0 \mod 2}} \EE(X_{i, j})
		= \begin{cases}
			\frac{n^2}{2} \cdot \frac{1}{16} & n \equiv 0 \mod 2 \\
			\frac{n^2 + 1}{2} \cdot \frac{1}{16} & n \equiv 0 \mod 2 \\
		\end{cases}
	\]
	וכן באופן דומה נובע
	\[
		\EE(S_1)
		= \begin{cases}
			\frac{n^2}{2} \cdot \frac{1}{16} & n \equiv 0 \mod 2 \\
			\frac{n^2 - 1}{2} \cdot \frac{1}{16} & n \equiv 0 \mod 2 \\
		\end{cases}
	\]
	אילו נניח ש־$n \equiv 0 \mod 2$ אז
	\begin{align*}
		\PP(X \ge \frac{n^2}{16} + cn)
		& \le \PP(S_0 - \frac{n^2}{32} \ge \frac{1}{2} cn) + \PP(S_1 \ge \frac{n^2}{32} + \frac{1}{2} cn) \\
		& = \PP(S_0 - \frac{n^2}{32} \ge \frac{1}{2} cn) + \PP(S_1 \ge \frac{n^2}{32} + \frac{1}{2} cn) \\
		& \le \exp(-\frac{\frac{1}{4} c^2 n^2}{2 \frac{n^2}{2}}) + \exp(-\frac{\frac{1}{4} c^2 n^2}{2 \frac{n^2}{2}}) \\
		& = 2\exp(-\frac{c^2}{4})
	\end{align*}
	כפי שרצינו להוכיח, לכן נותר לנו להניח ש־$n \equiv 1 \mod 2$, במקרה זה
	\begin{align*}
		\PP(X \ge \frac{n^2}{16} + cn)
		& \le \PP(S_0 \ge \frac{n^2 + 1}{32} + \frac{1}{2} cn + \frac{c}{2}) + \PP(S_1 \ge \frac{n^2 - 1}{32} + \frac{1}{2} cn - \frac{1}{2} c) \\
		& \le \exp(\frac{\frac{1}{4} c^2 {(n + 1)}^2}{2 \frac{n^2 + 1}{2}}) + \exp(\frac{\frac{1}{4} c^2 {(n - 1)}^2}{2 \frac{n^2 - 1}{2}}) \\
		& = \exp(-\frac{c^2 {(n + 1)}^2}{4(n^2 + 1)}) + \exp(-\frac{c^2 {(n - 1)}^2}{4(n^2 - 1)}) \\
		& = 2 \exp(- \frac{c^2}{4})
	\end{align*}
	כפי שרצינו.

	נבחין כי ישנה דרך נוספת להוכחת הטענה.
	נניח שהצלעות של הטבלה הן בלתי־תלויות, דהינו הטבלה מהצורה:
	\begin{center}
		\begin{tabular}{||c || c || c ||}
			\hline\hline
			A & B & C \\
			\hline\hline
			A & B & C \\
			\hline\hline
			A & B & C \\
			\hline\hline
		\end{tabular}
	\end{center}
	נבחין שמתקיים
	\[
		\EE(Y_{i, j}) = \frac{1}{4}
	\]
	ולכן
	\[
		\EE(X_{i, j})
		= \EE(Y_{i, j}) \cdot \EE(Y_{i + 1, j}) \cdot \EE(Y_{i, j + 1}) \cdot \EE(Y_{i + 1, j + 1})
		= \frac{1}{16}
	\]
	ובהתאם
	\[
		\EE(X)
		= \sum_{1 \le i, j \le n} X_{i, j}
		= \frac{n^2}{16}
	\]
	וכמובן גם מההגדרה לכל $1 \le i, j \le n$ גם $|X_{i, j}| \le 1$, אז תנאי אי־שוויון הופדינג חלים ומתקיים
	\[
		\PP(X - \frac{n^2}{16} \ge cn)
		= \PP(X \ge cn + \frac{n^2}{16})
		\le e^{-\frac{c^2 n^2}{2 n^2}}
		\le 2e^{-\frac{c^2}{4}}
	\]
	כפי שרצינו להראות.
\end{proof}

\question{}
שיכור עומד על ציר המספרים, בכל יום הוא צועד צעד אחד שמאלה בהסתברות $p$ או שני צעדים ימינה בהסתברות $1 - p$, באופן בלתי־תלוי בימים הקודמים. \\
נוכיח שההסתברות שהשיכור נמצא בראשית ביום ה־$3n$ דועכת מעריכית לכל $p \ne \frac{2}{3}$.
\begin{proof}
	נגדיר $X_n$ המשתנה המקרי המייצג כמה זז השיכור ביום ה־$n$, כלומר
	\[
		\PP(X_n = s)
		= \begin{cases}
			p & s = -1 \\
			1 - p & s = 2
		\end{cases}
	\]
	וכן נגדיר $Y_n = \sum_{i = 1}^{n} X_n$, סך המרחק שזז השיכור לאורך $n$ ימים.
	נבחין ש־$\frac{X_n + 1}{3} \sim Ber(1 - p)$, ולכן $\frac{Y_n + n}{3} \sim Bin(n, 1 - p)$.
	לכן נסיק
	\[
		\EE(\frac{Y_n + n}{3})
		= n (1 - p)
		\iff
		\EE(Y_n + n) = n (3 - 3p)
		\iff
		\EE(XY_n) = n(2 - 3p)
	\]
	ולכן $\EE(Y_n - n(2 - 3p)) = 0$, וכן $\frac{1}{3} |X_n - (2 - 3p)| \le 1$ כמעט תמיד, ולכן
	\[
		\PP(Y_n = 0)
		\le \PP(Y_n \ge 0)
		= \PP(\frac{1}{3}(Y_n - n(2 - 3p)) \ge \frac{1}{3} n(2 - 3p))
		\le \exp(-\frac{n^2 {(2 - 3p)}^2}{9 \cdot 2n})
		\xrightarrow[p \ne \frac{2}{3}]{n \to \infty} e^{-\infty}
		= 0
	\]
	ומצאנו שביום ה־$n$ הסיכוי שהשיכור נמצא בראשית דועך לאפס לכל $p \ne \frac{2}{3}$, לכן בפרט ביום ה־$3n$ הוא גם דועך.
\end{proof}

\end{document}
